% chapters/07/section-03.tex
\section{Menyelesaikan Tabrakan Sumbu demi Sumbu}
\label{sec:per-sumbu}

Inilah bagian inti bab ini. Penyelesaian tabrakan sumbu demi sumbu adalah cara
yang membuat benda dapat mendarat di atas bongkahan sekaligus terhalang oleh
sisinya.

\subsection{Persoalannya}

Bayangkan pemain bergerak ke kanan sementara bongkahan berada tepat di
hadapannya. Apabila pemeriksaan dilakukan terhadap seluruh perpindahan
sekaligus, program hanya mengetahui bahwa kedua persegi beririsan. Program
tidak mengetahui dari sisi mana pemain datang: apakah dari kiri sehingga
seharusnya dihalangi, atau dari atas sehingga seharusnya mendarat.

Akibatnya muncul dua kesalahan yang khas. Apabila program selalu meletakkan
pemain di atas bongkahan, pemain yang menabrak sisi bongkahan akan tiba-tiba
terangkat ke atasnya. Apabila program selalu menghalangi pemain pada sisinya,
pemain yang jatuh di atas bongkahan akan tertahan di udaranya.

\subsection{Cara Menyelesaikannya}

Penyelesaiannya adalah memisahkan perpindahan menjadi dua tahap yang terpisah,
masing-masing dengan pemeriksaan sendiri. Urutannya diperlihatkan pada
Gambar~\ref{fig:tabrakan-per-sumbu}.

\begin{figure}[htbp]
\centering
% chapters/07/figures/tabrakan-per-sumbu.tikz
% Penyelesaian tabrakan sumbu demi sumbu: gerak mendatar diperiksa lebih
% dahulu, baru gerak tegak.
\begin{tikzpicture}[
  bongkah/.style={rectangle, draw=ubLine, thick, fill=ubGray},
  pemain/.style={rectangle, draw=ubPrimary, thick, fill=ubPrimary!25},
  pemainlama/.style={rectangle, draw=ubPrimary, dashed, thin},
  panah/.style={-{Stealth[length=2.2mm]}, thick, draw=ubAccent},
  judul/.style={font=\footnotesize\bfseries, ubPrimary, align=center},
  label/.style={font=\scriptsize, align=left}
]

% ---------- Panel A: sumbu mendatar ----------
\begin{scope}
  \node[judul, anchor=south] at (2.4,2.15)
    {1. Sumbu mendatar diperiksa lebih dahulu};

  % Bongkahan
  \draw[bongkah] (3.2,0) rectangle (5.6,2.0);
  \node[label, anchor=north] at (4.4,-0.15) {bongkahan};

  % Posisi sebelum dan sesudah penyelesaian
  \draw[pemainlama] (2.55,0.55) rectangle (3.05,1.45);
  \draw[pemain] (2.7,0.55) rectangle (3.2,1.45);

  \draw[panah] (2.2,1.7) -- (3.0,1.7);
  \node[label, anchor=south west] at (2.25,1.75)
    {bergerak ke kanan};

  % Tanda sisi yang ditempelkan
  \draw[ubRed, very thick] (3.2,0.4) -- (3.2,1.6);
  \node[label, ubRed, anchor=north east] at (3.15,0.35)
    {badan ditempelkan\\pada sisi kiri bongkahan};
\end{scope}

% ---------- Panel B: sumbu tegak ----------
\begin{scope}[xshift=8.4cm]
  \node[judul, anchor=south] at (2.4,2.15)
    {2. Sumbu tegak diperiksa kemudian};

  \draw[bongkah] (0.6,0) rectangle (4.4,0.9);
  \node[label, anchor=north] at (2.5,-0.15) {bongkahan};

  % Posisi sebelum dan sesudah penyelesaian
  \draw[pemainlama] (1.75,0.9) rectangle (2.45,1.8);
  \draw[pemain] (1.75,0.45) rectangle (2.45,1.35);

  \draw[panah] (2.85,1.9) -- (2.85,1.3);
  \node[label, anchor=west] at (2.95,1.6) {jatuh};

  \draw[ubGreen, very thick] (0.6,0.9) -- (4.4,0.9);
  \node[label, ubGreen, anchor=south west] at (2.6,0.95)
    {kaki diletakkan\\pada permukaan};
\end{scope}
\end{tikzpicture}

\caption{Penyelesaian tabrakan sumbu demi sumbu}
\label{fig:tabrakan-per-sumbu}
\end{figure}

\begin{langkah}
  \item Pindahkan benda pada \textbf{sumbu mendatar} saja, lalu periksa
        tabrakan. Apabila terjadi, sisi tabrakan dapat dipastikan dari arah
        gerak: bergerak ke kanan berarti menabrak dari kiri, dan sebaliknya.
        Perbaiki posisinya pada sumbu itu saja.
  \item Setelah sumbu mendatar selesai dan tidak ada lagi tumpang tindih,
        pindahkan benda pada \textbf{sumbu tegak}. Periksa kembali, dan
        perbaiki pada sumbu tegak.
\end{langkah}

Kuncinya ada pada urutan: perpindahan pada satu sumbu diselesaikan seluruhnya
sebelum sumbu berikutnya disentuh. Dengan begitu, pada setiap pemeriksaan hanya
ada satu sumbu yang perlu dipertimbangkan, dan sisi tabrakan menjadi jelas.

\subsection{Penerapannya pada Kode}

\begin{kodekeluaran}
# --- Sumbu mendatar ---
self.position.x += self.velocity.x * delta_seconds
self.rect.centerx = round(self.position.x)

for platform in platforms:
    if self.rect.colliderect(platform.rect):
        if self.velocity.x > 0:
            self.rect.right = platform.rect.left      # menabrak dari kiri
        elif self.velocity.x < 0:
            self.rect.left = platform.rect.right      # menabrak dari kanan
        self.position.x = self.rect.centerx
\end{kodekeluaran}

Perbaikan posisi dilakukan pada \code{rect}, lalu nilainya
disalin kembali ke \code{position}. Posisi tetap menjadi sumber kebenaran yang
berupa pecahan, sedangkan \code{rect} menjadi alat untuk memeriksa dan
memperbaiki. Itulah pembagian tugas yang telah Anda pelajari pada
Bagian~5.4, kini dipakai dari arah sebaliknya.

\begin{kodekeluaran}
# --- Sumbu tegak ---
self.velocity.y += GRAVITY * delta_seconds
self.position.y += self.velocity.y * delta_seconds
self.rect.centery = round(self.position.y)

self.on_ground = False
for platform in platforms:
    if self.rect.colliderect(platform.rect):
        if self.velocity.y > 0:
            self.rect.bottom = platform.rect.top      # mendarat di atas
            self.velocity.y = 0.0
            self.on_ground = True
        elif self.velocity.y < 0:
            self.rect.top = platform.rect.bottom      # membentur langit-langit
            self.velocity.y = 0.0
        self.position.y = self.rect.centery
\end{kodekeluaran}

Tiga hal perlu diperhatikan pada bagian tegak.

\begin{enumerate}
  \item \textbf{Kecepatan tegak dinolkan ketika tabrakan diselesaikan.} Tanpa
        itu, benda akan terus tertekan ke dalam bongkahan pada bingkai
        berikutnya, dan posisinya diperbaiki berulang kali pada setiap bingkai.
        Program tetap tampak benar, tetapi benda bergetar halus di permukaan.
  \item \textbf{Penanda \code{on_ground} diatur di sini}, bukan dihitung
        terpisah. Cara itu lebih tepat daripada membandingkan posisi dengan
        permukaan bongkahan, sebab penandanya dihasilkan oleh pemeriksaan yang
        sesungguhnya.
  \item \textbf{Gravitasi ditambahkan sebelum gerak tegak.}
        Urutan itu tetap mengikuti kaidah Euler semi-implisit pada
        Bagian~2.4, dengan pemeriksaan tabrakan disisipkan sesudah perpindahan.
\end{enumerate}

\subsection{Memeriksa Kebenaran Penanganan secara Otomatis}

Penanganan tabrakan yang salah jarang menimbulkan pesan kesalahan. Karena itu,
buku ini menganjurkan satu kebiasaan yang sederhana tetapi sangat berguna:
\textbf{periksa sendiri kebenarannya di dalam program}, lalu tampilkan
hasilnya.

\begin{kodekeluaran}
def count_overlaps(player, platforms):
    return sum(1 for platform in platforms
               if player.rect.colliderect(platform.rect))
\end{kodekeluaran}

Fungsi itu menghitung berapa bongkahan yang badannya bertumpang tindih dengan
pemain. Nilainya seharusnya selalu nol: sesudah penyelesaian tabrakan, pemain
tidak boleh berada di dalam bongkahan mana pun. Apabila nilainya bukan nol pada
suatu bingkai, penanganan tabrakan Anda belum benar.

Kebiasaan semacam itu berguna jauh di luar bab ini. Sebuah program yang
memeriksa keadaan dirinya sendiri dan melaporkan penyimpangan akan menemukan
kesalahan lebih cepat daripada program yang hanya menampilkan gambar di layar.
Pada program bab ini, penghitung itu ditampilkan pada terminal setiap kali
program berakhir.
